\documentclass[11pt]{article}
% ======================================================================
%  examstyle.tex — EPFL-style exam layout for the weekly Time Series exams
%  Shared by every Exam_Week_NN(.tex / _Solutions.tex). Compiles with pdflatex.
% ======================================================================
\usepackage[utf8]{inputenc}
\usepackage[T1]{fontenc}
\usepackage[english]{babel}
\usepackage[a4paper,margin=2.4cm]{geometry}
\usepackage{amsmath,amssymb,amsthm,mathtools}
\usepackage{bm}
\usepackage{enumitem}
\usepackage{array}

\setlength{\parindent}{0pt}
\setlength{\parskip}{0.5em}
\emergencystretch=3em
\allowdisplaybreaks[1]

% ---------- math macros (same names as the rest of the project) ----------
\newcommand{\E}{\mathbb{E}}
\DeclareMathOperator{\Var}{Var}
\DeclareMathOperator{\Cov}{Cov}
\DeclareMathOperator{\Corr}{Corr}
\DeclareMathOperator*{\argmin}{arg\,min}
\DeclareMathOperator{\MSE}{MSE}
\DeclareMathOperator{\trace}{tr}
\newcommand{\Reals}{\mathbb{R}}
\newcommand{\Ints}{\mathbb{Z}}
\newcommand{\Comp}{\mathbb{C}}
\newcommand{\Nat}{\mathbb{N}}
\newcommand{\Prob}{\mathbb{P}}
\newcommand{\Tr}{^{\mathsf{T}}}
\newcommand{\conj}[1]{\overline{#1}}
\newcommand{\abs}[1]{\left\lvert #1 \right\rvert}
\newcommand{\norm}[1]{\left\lVert #1 \right\rVert}
\newcommand{\ind}{\mathbf{1}}
\newcommand{\eps}{\varepsilon}
\newcommand{\diff}{\nabla}
\newcommand{\acvs}{\gamma}
\newcommand{\acf}{\rho}
\newcommand{\pacf}{\alpha}
\newcommand{\sdf}{\mathcal{S}}
\newcommand{\Bsh}{B}
\newcommand{\iid}{\overset{\text{iid}}{\sim}}

\setlist[enumerate]{leftmargin=2em,topsep=2pt,itemsep=4pt}
\setlist[itemize]{leftmargin=1.6em,topsep=2pt,itemsep=2pt}

\newcounter{exo}

% ---------- exam cover header :  \examheader{exam title}{duration} ----------
\newcommand{\examheader}[2]{%
  \thispagestyle{empty}
  \begin{center}
    {\Large\bfseries Time Series}\\[3pt]
    Spring Semester 2026, EPFL\\
    \'Ecole Polytechnique F\'ed\'erale de Lausanne\\
    Prof.\ Dr.\ Sofia C.\ Olhede\\[7pt]
    \hrule height 0.8pt
    \vspace{9pt}
    {\large\bfseries Practice Exam --- #1}\\[4pt]
    {\normalsize Duration: #2 \quad$\cdot$\quad No calculator \quad$\cdot$\quad Answer on paper}\\
    \vspace{8pt}
    \hrule height 0.8pt
  \end{center}
  \vspace{14pt}
  \begin{tabular}{@{}p{8.2cm}@{}p{8cm}@{}}
    Name:\enspace\hrulefill & Surname:\enspace\hrulefill
  \end{tabular}\\[14pt]
  SCIPER (Student Card Number):\enspace\hrulefill
  \vspace{16pt}
}

% ---------- points table (fixed: 4 problems) ----------
\newcommand{\pointstable}{%
  \begin{center}
  \renewcommand{\arraystretch}{1.5}
  \begin{tabular}{|p{5cm}|p{4cm}|}
    \hline
    \textbf{Exercise} & \textbf{Points}\\\hline
    1 & \\\hline
    2 & \\\hline
    3 & \\\hline
    4 & \\\hline
    \textbf{Total} & \\\hline
  \end{tabular}
  \end{center}
}

% ---------- problem heading (exam: one problem per page) ----------
\newcommand{\problem}[1]{%
  \clearpage
  \stepcounter{exo}%
  \begin{center}\textbf{\large Problem \theexo\ (#1 points)}\end{center}
  \vspace{4pt}
}
% ---------- answer space: fill the statement page + one blank answer page ----------
% Put \answerpage at the END of each problem so every exercise gets its own page(s).
\newcommand{\answerpage}{\par\vfill\newpage\null\newpage}

% ---------- solutions document ----------
\newcommand{\solheader}[1]{%
  \begin{center}
    {\Large\bfseries Time Series --- Solutions}\\[3pt]
    {\large Practice Exam --- #1}\\[2pt]
    EPFL, MATH-342 \quad$\cdot$\quad Prof.\ S.\ C.\ Olhede\\[5pt]
    \hrule height 0.8pt
  \end{center}
  \vspace{10pt}
}
\newcommand{\solproblem}[1]{%
  \bigskip
  \stepcounter{exo}%
  {\large\bfseries Problem \theexo\ (#1 points)}\par\nobreak\vspace{2pt}
}
% Rappel de l'énoncé avant la solution (dans les corrigés)
\newcommand{\statementstart}{\par\vspace{1pt}{\bfseries\itshape Statement.}\par\nobreak\begingroup\small\itshape}
\newcommand{\statementend}{\par\endgroup\vspace{5pt}\hrule\vspace{6pt}{\bfseries Solution.}\par\nobreak\vspace{2pt}}

\begin{document}

\examheader{Week 3: Backshift, Wold Decomposition, Causality \& Invertibility}{60 minutes}
\pointstable
\begin{center}\itshape Justify every answer. Throughout, $\eps_t$ is a mean-zero white noise of variance $\sigma^2$.\end{center}

% ---------------------------------------------------------------
\problem{5}
\textbf{(a)} Define the \emph{backshift operator} $\Bsh$ and, for an $\mathrm{ARMA}(p,q)$ process written $\Phi(\Bsh)X_t=\Theta(\Bsh)\eps_t$, write down the characteristic polynomials $\Phi(z)$ and $\Theta(z)$. State the \emph{root criteria} (in terms of the unit circle) for (i) stationarity/causality and (ii) invertibility. \hfill(2 pts)

\textbf{(b)} Consider the $\mathrm{AR}(2)$ process
\[
X_t=\tfrac23 X_{t-1}-\tfrac16 X_{t-2}+\eps_t .
\]
Write its AR polynomial $\Phi(z)$, find the roots, and decide whether the process is stationary. \hfill(2 pts)

\textbf{(c)} Is the process of part (b) invertible? Justify in one sentence (no computation needed). \hfill(1 pt)
\answerpage

% ---------------------------------------------------------------
\problem{5}
For the $\mathrm{ARMA}(2,1)$ process
\[
(1-0.7\Bsh+0.1\Bsh^2)\,X_t=(1+0.5\Bsh)\,\eps_t,
\]

\textbf{(a)} factor the AR polynomial $\Phi(z)=1-0.7z+0.1z^2$, locate its roots, and confirm that the process is causal. \hfill(2 pts)

\textbf{(b)} Using a \emph{partial-fraction} expansion of the transfer function $G(z)=\Theta(z)/\Phi(z)$, find a closed form for the $\mathrm{MA}(\infty)$ weights $\psi_j$ in $X_t=\sum_{j\ge0}\psi_j\eps_{t-j}$, and report $\psi_0,\psi_1,\psi_2$. \hfill(3 pts)
\answerpage

% ---------------------------------------------------------------
\problem{5}
Let $\{X_t\}$ be the causal $\mathrm{AR}(2)$ process
\[
X_t=\tfrac56 X_{t-1}-\tfrac16 X_{t-2}+\eps_t .
\]

\textbf{(a)} Show that the autocovariance sequence satisfies the homogeneous recurrence
$\acvs_\tau=\tfrac56\,\acvs_{\tau-1}-\tfrac16\,\acvs_{\tau-2}$ for $\tau\ge1$, and obtain the boundary relation at $\tau=0$ involving $\sigma^2$. \hfill(2 pts)

\textbf{(b)} Solve for $\acvs_0,\acvs_1$ and hence give the closed-form expression of $\acvs_\tau$ for all $\tau\ge0$. \hfill(3 pts)

\emph{Hint: the right-hand side of the model equation only contributes through $\E[\eps_t X_t]=\sigma^2$ (causality).}
\answerpage

% ---------------------------------------------------------------
\problem{5}
\textbf{(Revision of Week 2 + this week.)} Let $X_t=\eps_t-\theta\eps_{t-1}$ be an $\mathrm{MA}(1)$ with $\theta=\tfrac12$.

\textbf{(a)} Compute the autocovariance sequence $\acvs_\tau$ for all $\tau$ and the autocorrelation $\acf_1$. \hfill(1.5 pts)

\textbf{(b)} Prove that \emph{any} $\mathrm{MA}(1)$ of the form $\eps_t-\theta\eps_{t-1}$ satisfies $\abs{\acf_1}\le\tfrac12$, and state qualitatively how the PACF of an $\mathrm{MA}(1)$ behaves (cut off or tail off?). \hfill(1.5 pts)

\textbf{(c)} Is $\{X_t\}$ invertible? Then exhibit a \emph{different} $\mathrm{MA}(1)$ process $Y_t=\eta_t-\vartheta\eta_{t-1}$ (with an appropriately chosen white-noise variance) that has \emph{exactly the same} autocovariance sequence as $\{X_t\}$, and explain which of the two is the invertible representation. \hfill(2 pts)
\answerpage

\end{document}
